\documentclass[12pt]{article}
\usepackage{amsmath, amssymb}
\usepackage{geometry}
\usepackage{booktabs}
\usepackage{array}
\usepackage{enumitem}
\usepackage{parskip}

\geometry{margin=1in}

\title{Econ 20200 --- Problem Set 1 Solutions}
\date{Due April 3rd}

\begin{document}
\maketitle

%% ============================================================
\section*{Problem 1 (5 points) --- Skilled vs.\ Unskilled Labor}
%% ============================================================

The production function is
\[
    Y = z(K + 80000\,L_u)^{0.6}\,L_s^{0.4}.
\]

%% ----------------------------
\subsection*{Part (a) --- Marginal Products}
%% ----------------------------

Taking partial derivatives:

\[
    MPK
        = \frac{\partial Y}{\partial K}
        = 0.6\,z\,(K+80000\,L_u)^{-0.4}\,L_s^{0.4}
\]

\[
    MP_{L_u}
        = \frac{\partial Y}{\partial L_u}
        = 0.6\,z\cdot 80000\cdot(K+80000\,L_u)^{-0.4}\,L_s^{0.4}
        = 80000 \cdot MPK
\]

\[
    MP_{L_s}
        = \frac{\partial Y}{\partial L_s}
        = 0.4\,z\,(K+80000\,L_u)^{0.6}\,L_s^{-0.6}
\]

Note the structural relationship $MP_{L_u} = 80000 \cdot MPK$, which reflects
the fact that capital and unskilled labor are \emph{perfect substitutes}
(each unskilled worker is equivalent to 80{,}000 units of capital in this
production function). By contrast, $MP_{L_s}$ is \emph{increasing} in $K$,
confirming that capital and skilled labor are \emph{complements}.

%% ----------------------------
\subsection*{Part (b) --- Numerical Comparison, $L_s = 0.6$ in Both Periods}
%% ----------------------------

\textbf{1990:} $z=50,\; K=25000,\; L_u=1,\; L_s=0.6$.

\[
    K + 80000\,L_u = 25000 + 80000 = 105000
\]

Key computed values:
$(105000)^{0.4} \approx 102.0$,\;
$(0.6)^{0.4} \approx 0.8152$,\;
$(105000)^{0.6} \approx 1030.4$,\;
$(0.6)^{-0.6} \approx 1.3587$.

\begin{align*}
    MP_{L_u}^{1990}
        &= 50 \times 0.6 \times 80000 \times (105000)^{-0.4} \times (0.6)^{0.4} \\
        &= 50 \times 0.6 \times 80000 \times 0.00980 \times 0.8152
         \approx 19{,}180 \\[4pt]
    MP_{L_s}^{1990}
        &= 50 \times 0.4 \times (105000)^{0.6} \times (0.6)^{-0.6} \\
        &= 50 \times 0.4 \times 1030.4 \times 1.3587
         \approx 28{,}000
\end{align*}

\textbf{2020:} $z=60,\; K=50000,\; L_u=1,\; L_s=0.6$.

\[
    K + 80000\,L_u = 50000 + 80000 = 130000
\]

Key computed values:
$(130000)^{-0.4} \approx 0.008984$,\;
$(130000)^{0.6} \approx 1167.6$.

\begin{align*}
    MP_{L_u}^{2020}
        &= 60 \times 0.6 \times 80000 \times 0.008984 \times 0.8152
         \approx 21{,}090 \\[4pt]
    MP_{L_s}^{2020}
        &= 60 \times 0.4 \times 1167.6 \times 1.3587
         \approx 38{,}090
\end{align*}

\begin{center}
\begin{tabular}{lccc}
\toprule
 & $MP_{L_u}$ & $MP_{L_s}$ & Skill Premium ($MP_{L_s}/MP_{L_u}$) \\
\midrule
1990 & 19,180 & 28,000 & 1.46 \\
2020 & 21,090 & 38,090 & 1.81 \\
\% Change & $+10\%$ & $+36\%$ & $\uparrow$ \\
\bottomrule
\end{tabular}
\end{center}

\textbf{Comment:} Both capital doubling and TFP growth raised the marginal
products of all workers, but $MP_{L_s}$ rose far more than $MP_{L_u}$ because
capital \emph{complements} skilled labor while \emph{substituting} for
unskilled labor. As a result, the skill premium widened from 1.46 to 1.81,
mirroring the real-world rise in wage inequality over 1990--2020.

%% ----------------------------
\subsection*{Part (c) --- 2020 with $L_s = 0.7$}
%% ----------------------------

Now $L_s^{2020} = 0.7$. Updated key values:
$(0.7)^{0.4} \approx 0.8669$,\; $(0.7)^{-0.6} \approx 1.2388$.

\begin{align*}
    MP_{L_u}^{2020}
        &= 60 \times 0.6 \times 80000 \times 0.008984 \times 0.8669
         \approx 22{,}430 \\[4pt]
    MP_{L_s}^{2020}
        &= 60 \times 0.4 \times 1167.6 \times 1.2388
         \approx 34{,}710
\end{align*}

\begin{center}
\begin{tabular}{lccc}
\toprule
 & $MP_{L_u}$ & $MP_{L_s}$ & Skill Premium \\
\midrule
1990 & 19,180 & 28,000 & 1.46 \\
2020 (part c) & 22,430 & 34,710 & 1.55 \\
\% Change & $+17\%$ & $+24\%$ & $\uparrow$ (less) \\
\bottomrule
\end{tabular}
\end{center}

\textbf{Comparison with Part (b):} When the supply of skilled workers grows
from 0.6 to 0.7, $MP_{L_u}$ rises \emph{more} (higher $L_s$ amplifies the
productivity of the whole system, which benefits unskilled workers too), while
$MP_{L_s}$ rises \emph{less} due to diminishing marginal returns. The skill
premium increases only to 1.55 versus 1.81 in part (b). If the supply of
skilled workers keeps pace with demand from capital accumulation, wage
inequality is moderated.

%% ----------------------------
\subsection*{Part (d) --- Effect of TFP Growth ($z \uparrow$)}
%% ----------------------------

From the expressions in part (a), $z$ enters multiplicatively in \emph{both}
marginal products:
\[
    MP_{L_u} \propto z, \qquad MP_{L_s} \propto z.
\]

An increase in TFP raises the marginal product of \emph{both} types of labor
proportionally, shifting labor demand outward for each type. Because $z$
cancels in the skill premium ratio $MP_{L_s}/MP_{L_u}$, TFP growth alone does
\emph{not} change the relative demand for skilled versus unskilled labor ---
it is neutral between skill types. Firms hire more of both workers.

%% ----------------------------
\subsection*{Part (e) --- Effect of Capital Accumulation ($K \uparrow$)}
%% ----------------------------

\begin{itemize}[leftmargin=*]
    \item $MP_{L_u} \propto (K + 80000\,L_u)^{-0.4}$: as $K$ rises, this
    term \emph{decreases}, so $MP_{L_u}$ \emph{falls}. Capital substitutes
    for unskilled labor --- more machines reduce the marginal value of
    additional unskilled workers.

    \item $MP_{L_s} \propto (K + 80000\,L_u)^{0.6}$: as $K$ rises, this
    term \emph{increases}, so $MP_{L_s}$ \emph{rises}. Capital complements
    skilled labor --- more machines make trained operators more productive.
\end{itemize}

Capital accumulation raises the demand for skilled labor and lowers the demand
for unskilled labor, widening the skill premium.

%% ============================================================
\section*{Problem 2 (6 points) --- Firm Optimization: $Y = 100L^{0.6}$, $w = 10$}
%% ============================================================

%% ----------------------------
\subsection*{Part (a) --- Baseline}
%% ----------------------------

The firm's problem is:
\[
    \max_{L} \; \pi(L) = 100L^{0.6} - 10L
\]

\textbf{First-order condition:}
\[
    \frac{d\pi}{dL} = 60\,L^{-0.4} - 10 = 0
    \;\Longrightarrow\;
    L^{0.4} = 6
    \;\Longrightarrow\;
    L^* = 6^{2.5} = 36\sqrt{6} \approx 88.18
\]

\textbf{Output:}
\[
    Y^* = 100\,(6^{2.5})^{0.6} = 100 \cdot 6^{1.5} = 600\sqrt{6} \approx 1469.7
\]

\textbf{Maximized profit:}
\[
    \pi^* = 600\sqrt{6} - 10\cdot 36\sqrt{6} = 240\sqrt{6} \approx 587.9
\]

%% ----------------------------
\subsection*{Part (b) --- Government Subsidy of 0.5 per Worker}
%% ----------------------------

The government pays the firm 0.5 units per worker hired, so the effective
labor cost is $w_{\text{eff}} = 10 - 0.5 = 9.5$. The new problem is:
\[
    \max_{L} \; 100L^{0.6} - 9.5\,L
\]

\textbf{FOC:}
\[
    60\,L^{-0.4} = 9.5
    \;\Longrightarrow\;
    L^{0.4} = \frac{60}{9.5} \approx 6.316
    \;\Longrightarrow\;
    L^* \approx 100.1
\]

Since the effective wage fell ($9.5 < 10$), the firm hires \textbf{more}
labor, produces \textbf{more} output, and earns \textbf{higher} profit
relative to the baseline.

%% ----------------------------
\subsection*{Part (c) --- Environmental Compliance}
%% ----------------------------

Each production worker requires an additional 0.1 cleanup workers, so the
effective cost per production worker is $(1+0.1)\,w = 11$. Output is still
$Y = 100L^{0.6}$ where $L$ is production workers. The firm solves:
\[
    \max_{L} \; 100L^{0.6} - 11\,L
\]

\textbf{FOC:}
\[
    60\,L^{-0.4} = 11
    \;\Longrightarrow\;
    L^{0.4} = \frac{60}{11} \approx 5.455
    \;\Longrightarrow\;
    L^* \approx 69.4
\]

Since the effective cost rose ($11 > 10$), the firm hires \textbf{fewer}
production workers, produces \textbf{less} output, and earns \textbf{lower}
profit. The environmental regulation acts as a tax on labor, reducing
employment on the production line.

%% ============================================================
\section*{Problem 3 (6 points) --- Corner Solutions}
%% ============================================================

%% ----------------------------
\subsection*{Part (a) --- Dan: $u(c,h) = c + 0.8h$, $\pi = 0.8$}
%% ----------------------------

Dan's optimization problem is:
\[
    \max_{c,\,h} \; c + 0.8h
    \qquad \text{s.t.} \quad c \leq w(1-h) + 0.8
\]

\textbf{Lagrangian:}
\[
    \mathcal{L} = c + 0.8h + \lambda\bigl[w(1-h) + 0.8 - c\bigr]
\]

\textbf{First-order conditions:}
\begin{align*}
    \frac{\partial \mathcal{L}}{\partial c} &= 1 - \lambda = 0
        &&\Longrightarrow\; \lambda = 1 \\[4pt]
    \frac{\partial \mathcal{L}}{\partial h} &= 0.8 - \lambda w = 0
        &&\Longrightarrow\; w = 0.8
\end{align*}

The FOC pins down a specific wage ($w = 0.8$) rather than an interior choice
of $(c,h)$; for any other wage there is no interior solution. The
\textbf{second-order sufficient condition fails} because $u$ is linear ---
all second derivatives of the Lagrangian are identically zero, so the Hessian
is zero and we cannot confirm a strict interior maximum.

%% ----------------------------
\subsection*{Part (b) --- Dan's Labor Supply Curve}
%% ----------------------------

Since $u$ is linear, Dan compares his $MRS_{h,c} = 0.8$ to the wage $w$:

\begin{itemize}[leftmargin=*]
    \item If $w > 0.8$: leisure is ``too expensive''; Dan works all available
    time: $1-h = 1$.
    \item If $w = 0.8$: Dan is exactly indifferent; any $1-h \in [0,1]$ is
    optimal.
    \item If $w < 0.8$: leisure is relatively cheap; Dan takes all leisure:
    $1-h = 0$.
\end{itemize}

Dan's labor supply curve is a step function: zero hours for $w < 0.8$, the
full time endowment for $w > 0.8$, and any quantity in between at $w = 0.8$.

%% ----------------------------
\subsection*{Part (c) --- Nate: $u(c,h) = \min\{c,\,5h\}$, $\pi = 0$}
%% ----------------------------

Nate's indifference curves are \textbf{L-shaped} (perfect complements) with
kinks along the ray $c = 5h$.

\textbf{Argument that the optimum satisfies $c = 5h$:}
\begin{itemize}[leftmargin=*]
    \item If $c > 5h$: then $u = 5h$. Nate could reduce $c$ without
    reducing utility, saving income --- not optimal.
    \item If $c < 5h$: then $u = c$. Nate could reduce $h$ (work more),
    increase income and spend it on $c$, strictly raising utility ---
    not optimal.
\end{itemize}
Therefore the optimum must lie at the kink: $c = 5h$.

%% ----------------------------
\subsection*{Part (d) --- Nate's Labor Supply Curve}
%% ----------------------------

At the optimum, $c = 5h$ and the budget constraint binds ($\pi = 0$):
\[
    5h = w(1-h)
    \;\Longrightarrow\;
    5h = w - wh
    \;\Longrightarrow\;
    h^*(5 + w) = w
    \;\Longrightarrow\;
    h^* = \frac{w}{5+w}
\]
\[
    \boxed{1 - h^* = \frac{5}{5+w}}
\]

As $w$ increases, $5/(5+w)$ \textbf{decreases}, so Nate works \emph{less}
as the wage rises. His labor supply is downward-sloping from the outset.

\textbf{Income and substitution effects:} With perfect complements there is
\textbf{no substitution effect} --- the kink in the indifference curve means
relative price changes do not induce any reallocation along a given
indifference curve. The \textbf{income effect} therefore dominates completely:
a higher wage makes Nate richer, so he ``buys'' more leisure and supplies
fewer hours of labor.

%% ============================================================
\section*{Problem 4 (8 points) --- Jane: $u(c,h) = \sqrt{c} + 4\sqrt{h}$}
%% ============================================================

Jane's budget constraint is $c \leq w(1-h) + \pi$.

%% ----------------------------
\subsection*{Part (a) --- $w = 20$, $\pi = 5$}
%% ----------------------------

\textbf{Interior optimality condition} ($MRS_{h,c} = w$):
\[
    MRS_{h,c}
        = \frac{MU_h}{MU_c}
        = \frac{2/\sqrt{h}}{1/(2\sqrt{c})}
        = \frac{4\sqrt{c}}{\sqrt{h}}
        = 20
    \;\Longrightarrow\;
    \sqrt{c} = 5\sqrt{h}
    \;\Longrightarrow\;
    c = 25h
\]

\textbf{Budget constraint (binding):}
\[
    c = 20(1-h) + 5 = 25 - 20h
\]

Substituting $c = 25h$:
\[
    25h = 25 - 20h
    \;\Longrightarrow\;
    45h = 25
    \;\Longrightarrow\;
    h^* = \frac{5}{9}
\]
\[
    \boxed{1 - h^* = \frac{4}{9} \approx 0.444,
    \qquad c^* = \frac{125}{9} \approx 13.89}
\]

%% ----------------------------
\subsection*{Part (b) --- Labor Supply Curve (general $w$, $\pi = 5$)}
%% ----------------------------

From the optimality condition $MRS_{h,c} = w$:
\[
    \frac{4\sqrt{c}}{\sqrt{h}} = w
    \;\Longrightarrow\;
    c = \frac{w^2}{16}\,h
\]

Substituting into the binding budget constraint $c = w(1-h) + 5$:
\[
    \frac{w^2}{16}\,h + wh = w + 5
    \;\Longrightarrow\;
    h\cdot\frac{w(w+16)}{16} = w + 5
    \;\Longrightarrow\;
    h^* = \frac{16(w+5)}{w(w+16)}
\]

\[
    \boxed{1 - h^* = \frac{w^2 - 80}{w(w+16)}}
\]

Jane supplies positive hours only when $w^2 > 80$, i.e., $w > 4\sqrt{5} \approx 8.94$.

\begin{center}
\begin{tabular}{cc}
\toprule
Wage $w$ & Hours worked $1-h^*$ \\
\midrule
10 & $20/260 \approx 0.077$ \\
12 & $64/336 \approx 0.190$ \\
16 & $176/512 \approx 0.344$ \\
20 & $320/720 \approx 0.444$ \\
24 & $496/960 \approx 0.517$ \\
\bottomrule
\end{tabular}
\end{center}

The labor supply curve is \textbf{upward-sloping}: the substitution effect
(higher wages make work more attractive relative to leisure) dominates the
income effect for this utility function.

%% ----------------------------
\subsection*{Part (c) --- Effect of Higher Non-Wage Income $\pi$}
%% ----------------------------

Repeating the derivation for general $\pi$, the labor supply becomes:
\[
    1 - h^* = \frac{w^2 - 16\pi}{w(w+16)}
\]

An increase in $\pi$ \textbf{reduces} $1-h^*$: Jane supplies fewer hours at
every wage level. This is a \textbf{pure income effect} --- with no change in
the wage, there is no substitution effect, and the higher non-wage income
allows Jane to afford more leisure. The labor supply curve \textbf{shifts
inward (left)}: Jane works less at every wage.

\end{document}